\documentclass[11pt,a4paper]{article}
\usepackage[utf8]{inputenc}
\usepackage[italian]{babel}
\usepackage{amsmath}
\usepackage{amsfonts}
\usepackage{amssymb}
\usepackage{graphicx}
\usepackage{geometry}
\usepackage{booktabs}
\usepackage{listings}
\usepackage{xcolor}
\usepackage{hyperref}


% Configurazione pulita e sicura per il codice C++
\definecolor{commentcolor}{rgb}{0.12, 0.53, 0.53}
\definecolor{keywordcolor}{rgb}{0.0, 0.0, 0.8}
\definecolor{stringcolor}{rgb}{0.63, 0.08, 0.08}
\definecolor{backgroundcolor}{rgb}{0.97, 0.97, 0.97}

\lstset{
    language=C++,
    backgroundcolor=\color{backgroundcolor},
    basicstyle=\ttfamily\small,
    keywordstyle=\color{keywordcolor}\bfseries,
    commentstyle=\color{commentcolor}\itshape,
    stringstyle=\color{stringcolor},
    numbers=left,
    numberstyle=\tiny\color{gray},
    stepnumber=1,
    numbersep=8pt,
    frame=single,
    breaklines=true,
    tabsize=2,
    showstringspaces=false,
    extendedchars=true,
    literate={->*}{{\texttt{->*}}}3 {.*}{{\texttt{.*}}}2
}


\geometry{margin=2.5cm}

% FIX: Rimossa la graffa in eccesso che rompeva il preambolo
\title{\textbf{Report Tecnico: Architettura Real-Time Dual-Arm Control}\\ \large Ottimizzazione FSM tramite Funzioni Membro e Analisi di Complessità}
\author{Sviluppo Software di Controllo Robotico}
\date{Giugno 2026}

\begin{document}

\maketitle
\tableofcontents
\newpage

% =========================================================================
% SEZIONE 1: INTRODUZIONE
% =========================================================================
\section{Introduzione}
Nei sistemi di controllo robotico operanti in ambiente deterministico e real-time (con frequenze operative comprese tra 200 Hz e 1 kHz), la minimizzazione della latenza computazionale e l'eliminazione del jitter di memoria sono requisiti fondamentali.

Questo documento analizza l'evoluzione dell'architettura software del modulo di controllo robotico denominato \texttt{DualArmControl}, sviluppato sul framework \texttt{mc\_rtc}. L'architettura è stata ottimizzata con un pattern basato su un \textbf{Doppio Puntatore a Funzione Membro} a esecuzione deterministica pre-allocata.


% =========================================================================
% SEZIONE 2: ARCHITETTURA DEL CODICE C++
% =========================================================================
\section{Architettura e Struttura del Codice C++}
La vecchia implementazione software si appoggiava su flag booleani locali per discriminare la prima esecuzione di uno stato dall'esecuzione ciclica. Il nuovo design introduce una separazione netta e strutturale tra la fase di ingresso \textit{one-shot} in uno stato (\textit{Entry}) e la successiva esecuzione ciclica (\textit{Run}).

\subsection{Pattern FSM a Doppio Puntatore}
L'astrazione si basa sull'alias di tipo \texttt{StateMethod}, definito come puntatore a una funzione membro della classe di controllo:

\begin{lstlisting}
using StateMethod = void (DualArmControl::*)();
StateMethod currentState;
\end{lstlisting}

Il meccanismo di transizione gestisce la sequenza logica di switch in modo sicuro:

\begin{lstlisting}
void DualArmControl::transitionTo(StateMethod entryMethod, 
                                  StateMethod runMethod) {
    if (!entryMethod || !runMethod) return;
    stateTimer = 0.0; 
    (this->*entryMethod)(); 
    currentState = runMethod; 
}
\end{lstlisting}

Grazie a questa astrazione, il loop principale di controllo \texttt{run()} viene liberato da qualsiasi ramificazione condizionale o istruzione di salto indotto:

\begin{lstlisting}
bool DualArmControl::run() {
    (this->*currentState)(); 
    return mc_control::MCController::run();
}
\end{lstlisting}

\subsection{Pre-allocazione e Azzeramento del Real-Time Jitter}
All'interno del costruttore e del metodo di \texttt{reset()}, tutti i task del solutore quadratico (QP) come \texttt{PostureTask}, \texttt{EndEffectorTask} e \texttt{ImpedanceTask} vengono allocati preventivamente nello spazio di memoria. 

Durante l'esecuzione sincrona nei singoli stati della FSM, il controllore si limita ad attivare o disattivare i task allocati inserendoli o rimuovendoli dal solutore di ottimizzazione attraverso i metodi nativi del framework:
\begin{lstlisting}
solver().addTask(leftEdgeTask);
solver().removeTask(leftEdgeTask);
\end{lstlisting}
Questo approccio previene allocazioni dinamiche a runtime.



% =========================================================================
% SEZIONE 4: DOCUMENTAZIONE DEL CONTROLLORE
% =========================================================================
\section{Documentazione del Controllore}
Questa sezione descrive le specifiche operative per l'utilizzo e l'estensione del modulo software.

\subsection{Flusso di Esecuzione della FSM}
La macchina a stati si sviluppa secondo tre macro-fasi sequenziali stabili:
\begin{enumerate}
    \item \textbf{IDLE}: Stato di sicurezza iniziale. Vincola il robot nella configurazione iniziale per 3.0 secondi.
    \item \textbf{INDEPENDENT}: Attiva i task operativi nello spazio cartesiano per gli end-effector destro e sinistro, guidando i bracci robotici verso coordinate operative indipendenti dello spazio di lavoro.
    \item \textbf{COLLABORATIVE}: Disattiva i vincoli cinematici indipendenti e inizializza i task di controllo ad impedenza. Definisce un sistema di riferimento virtuale centrale interpolando la traiettoria per la manipolazione cooperativa di un corpo rigido.
\end{enumerate}

\subsection{Guida per l'Aggiunta di un Nuovo Stato}
Per inserire un ulteriore stato logico all'interno della FSM (ad esempio uno stato denominato \texttt{stateEmergency}), è sufficiente seguire questa procedura in tre passaggi:

\subsubsection*{1. Definizione dell'Interfaccia nell'Header}
Includere i prototipi dei metodi di Entry e di Run all'interno dell'area privata del file di testata:
\begin{lstlisting}
void entryStateEmergency();
void stateEmergency();
\end{lstlisting}

\subsubsection*{2. Sviluppo della Logica nel Sorgente}
Implementare le funzioni isolando la logica di configurazione iniziale dai cicli di calcolo iterativi:
\begin{lstlisting}
void DualArmControl::entryStateEmergency() {
    solver().removeTask(leftImpedanceTask);
    solver().removeTask(rightImpedanceTask);
}

void DualArmControl::stateEmergency() {
    stateTimer += timeStep;
}
\end{lstlisting}

\subsubsection*{3. Chiamata della Routine di Transizione}
Eseguire lo switch dello stato chiamando il metodo di transizione preposto da un qualsiasi punto della FSM attiva utilizzando i relativi riferimenti di istanza:
\begin{lstlisting}
this->transitionTo(&DualArmControl::entryStateEmergency, 
                   &DualArmControl::stateEmergency);
\end{lstlisting}

\end{document}